\section{Mechanische Festkörpereigenschaften}
TODO


\section{Thermische Festkörpereigenschaften}
\begin{emphbox}
    \begin{description}
        \item[Anzahl an Atomen] $N = N_An$
        \item[Gaskonstante] $R$
        \item[Kelvin-Boltzmann-Konstante] $k_B$
        \item[Avogadro-Konstante] $N_A$
        \item[Stoffmenge in mol] $n$
    \end{description}
\end{emphbox}

\begin{sectionbox}
    \subsection{Spezifische Wärme(kapazität)}

    \subsubsection{Wärmekapazität}
    Erforderliche Wärmemenge (innere Energie), um die Temperatur eines Festkörpers um $1 K$ zu erhöhen
    \begin{equation*}
        C=\frac{\Delta U}{\Delta T}
    \end{equation*}

    \subsubsection{spezifische Wärmekapazität}
    \begin{equation*}
        c=\frac{1}{m} C
    \end{equation*}

    \subsubsection{Molwärme - Dulong-Petit-Gesetz}
    \begin{equation*}
        c_m = \frac{1}{n}C=3N_ak_B=3R=24,9\frac{\text{J}}{\text{mol K}}
    \end{equation*}
    Für hohe Temperaturen gilt diese Näherung (dass $c_m=$ const.). Für niedrigere Temperaturen muss man quantenmechanische Phänomene betrachten.

    \subsubsection{Debye-Näherung}
    sagt die Temperaturabhängigkeit der Wärmekapazität sowohl im Niedrig- als auch im Hochtemperaturlimes korrekt voraus.
    \begin{equation*}
        k_B\Theta = \hbar \omega_D
    \end{equation*}
    mit $\Theta$: materialspezifischer Parameter (Debye-Temperatur) und \\
    $\omega_D$: Debyesche Grenzfrequenz $\approx$ Frequenz der optischen Phononen
\end{sectionbox}

\begin{sectionbox}
    \subsection{Wärmeleitfähigkeit}
    Ein Material besitzt genau dann eine hohe Wärmeleitfähigkeit, wenn es eine hohe spezifische Wärme und eine geringe Streuung an Defekten oder anderen Phononen aufweist.
    \begin{equation*}
        \lambda = C\frac{v^2}{3}\tau = \frac{1}{3}Cvl
    \end{equation*}
    mit $\lambda$: Wärmeleitfähigkeit pro Volumen, $l$: mittlere freie Weglänge, $v$: Geschwindigkeit der Teilchen im Festkörper, $\tau$: Streuzeit
\end{sectionbox}


\section{Elektronen im Festkörper}
